\subsection{实现压缩、局部化圆商与 Lee--Yang 五片覆盖的最小三分裂}\label{subsec:conclusion-phase-ledger-branch-trisplitting}
沿用命题 \ref{prop:killo-implementation-vs-structural-half-circle-dimension}、
定理 \ref{thm:killo-godel-compression-not-finite-rank-homologizable}、
\ref{thm:conclusion-finite-prime-solenoid-terminal-object}、
\ref{thm:conclusion-finite-prime-localization-pontryagin-rigidity}、
\ref{thm:killo-godel-semidir-affine-2adic}、
\ref{thm:xi-time-part62d-leyang-branch-inverse-limit-five-torsors}
与推论 \ref{cor:killo-2adic-tate-hypercube-limit} 的记号。前述结果已经分别把实现层单射压缩、有限生成加法账本的同态化障碍、有限素数局部化的单圆连续商、Gödel 半直积在 \(T_2(E_{\min})\) 上的 \(2\)-进仿射表示，以及 Lee--Yang 分支塔的五片 \(2\)-进覆盖固定为闭式。把这些接口并读后，还可进一步抽出如下单链结论。

\begin{theorem}[实现压缩与同态线性化的严格分叉]\label{thm:conclusion-implementation-compression-hom-linearization-split}
在实现层与结构层半圆维口径下，有精确等式
\[
\cdim^{\mathrm{hom}}_{+}(\NN,+)=\frac12,
\qquad
\cdim^{\mathrm{hom}}_{+}(\NN_{\times})=+\infty,
\qquad
\cdim^{\mathrm{impl}}(\NN_{\times})=\frac12.
\]
因此，“可被单寄存器实现”与“可被有限加法账本忠实线性化”是两个彼此独立的性质：前者只测集合论单射，后者测乘法结构能否被压入有限生成的交换可消去加法对象。特别地，不存在任何有限生成交换可消去加法幺半群可作为 \(\NN_{\times}\) 的忠实对数账本。
\end{theorem}
\begin{proof}
前两条等式与最后一条等式正是命题 \ref{prop:killo-implementation-vs-structural-half-circle-dimension} 的内容。另一方面，定理 \ref{thm:killo-prime-freedom-non-finitizability} 已排除
\[
\NN_{\times}\hookrightarrow M
\]
对任意有限生成交换可消去幺半群 \(M\) 的单射幺半群同态，而定理 \ref{thm:killo-godel-compression-not-finite-rank-homologizable} 进一步排除了所有有限生成 \(\ZZ\oplus L\) 型加法账本上的忠实同态化。故实现层单寄存器可编址与结构层有限账本可线性化之间存在严格分叉。证毕。
\end{proof}

\begin{theorem}[有限素局部化的一圆商与全不连通核分解]\label{thm:conclusion-finite-localization-circlequotient-profinite-kernel}
设 \(S\subset \mathbb P\) 为有限素数集，并记
\[
G_S:=\ZZ[S^{-1}],
\qquad
\Sigma_S:=\widehat{G_S},
\qquad
P_S:=\prod_{p\in S}\ZZ_p.
\]
则存在短正合列
\[
0\longrightarrow P_S\longrightarrow \Sigma_S\longrightarrow \TT\longrightarrow 0,
\qquad
\cdim_{\mathrm{fr}}(G_S)=1.
\]
并且对任意两个有限素数集 \(S,T\subset \mathbb P\)，都有
\[
\Sigma_S/P_S\cong \TT \cong \Sigma_T/P_T,
\]
但同时
\[
\Sigma_S\cong \Sigma_T
\iff
S=T.
\]
换言之，所有有限素局部化共享同一个一维连续相位商，而全部算术特异性都被压入全不连通的 \(p\)-进核 \(P_S\) 中。
\end{theorem}
\begin{proof}
短正合列与 \(\cdim_{\mathrm{fr}}(G_S)=1\) 由定理 \ref{thm:conclusion-finite-prime-solenoid-terminal-object} 给出；其中特别说明 \(\Sigma_S\) 的任意非平凡连通紧 Lie 商都同构于 \(\TT\)。因此
\[
\Sigma_S/P_S\cong \TT
\]
对每个有限 \(S\) 都成立。另一方面，定理 \ref{thm:conclusion-finite-prime-localization-pontryagin-rigidity} 断言
\[
\Sigma_S\cong \Sigma_T\iff S=T,
\]
故连续圆商并不完备，真正完成分类的全部信息都滞留在全不连通核的 \(p\)-主因子中。证毕。
\end{proof}

\begin{corollary}[局部化数系的圆维坍缩与素支撑完备]\label{cor:conclusion-finite-localization-circle-dimension-collapse-profinite-support-completeness}
对任意有限素数集 \(S\subset \mathbb P\)，加法群
\[
G_S=\ZZ[S^{-1}]
\]
在有限秩扩展口径下都满足同一不变量
\[
\cdim_{\mathrm{fr}}(G_S)=1.
\]
然而其 Pontryagin 对偶中的全不连通核
\[
P_S=\prod_{p\in S}\ZZ_p
\]
却完全恢复素支撑：
\[
P_S\cong P_T
\iff
S=T.
\]
因此，连通圆因子把全部有限素局部化压缩到同一“圆维 \(1\)”类，而 profinite 核则把全部坏素支撑逐点记回。
\end{corollary}
\begin{proof}
定理 \ref{thm:conclusion-finite-localization-circlequotient-profinite-kernel} 已经同时给出
\[
\cdim_{\mathrm{fr}}(G_S)=1
\]
以及
\[
\Sigma_S\cong \Sigma_T\iff S=T.
\]
而 \(P_S=\prod_{p\in S}\ZZ_p\) 的 \(p\)-主因子非平凡当且仅当 \(p\in S\)，故 \(P_S\) 自身已经唯一决定 \(S\)。证毕。
\end{proof}

\begin{theorem}[Gödel 半直积的 \(2\)-进仿射化]\label{thm:conclusion-godel-semidirect-twoadic-affineization}
取相容基 \((P_n,Q_n)_{n\ge 1}\) 使
\[
P_n=2P_{n+1},
\qquad
Q_n=2Q_{n+1}.
\]
则有逆极限识别
\[
\varprojlim_n E_{\min}[2^n]\cong T_2(E_{\min})\cong \ZZ_2^2,
\]
且乘二映射 \([2]\) 在坐标上恰等于最高位截断。另一方面，若
\[
H:=\mathcal R\rtimes_{\Sigma}\NN
\]
为 prime-register 半直积，\(H_{\mathrm{can}}\subset H\) 为规范有限时深子单子，并取
\[
B=2^L>M,
\]
则映射
\[
\Psi(r,k)(x):=B^k x+\operatorname{ev}_B(r)
\]
给出忠实单子嵌入
\[
\Psi:H_{\mathrm{can}}\hookrightarrow \mathrm{Aff}\!\bigl(T_2(E_{\min})\bigr).
\]
因此，prime-register 侧的尺度推进与事件写入，正对应于同一 \(2\)-进 Tate 模块上的仿射膨胀与平移。
\end{theorem}
\begin{proof}
逆极限识别与“乘二即最高位截断”正是推论 \ref{cor:killo-2adic-tate-hypercube-limit} 的内容。规范子单子在 \(T_2(E_{\min})\) 上的忠实仿射表示，则由定理 \ref{thm:killo-godel-semidir-affine-2adic} 直接给出；其中阈值条件正是 \(B>M\)。故 Gödel 半直积与分辨率塔并非仅为类比，而是同一 \(2\)-进仿射动力学在两种坐标系中的严格实现。证毕。
\end{proof}

\begin{theorem}[Lee--Yang 分支塔的五片 dyadic 覆盖]\label{thm:conclusion-leyang-five-sheet-dyadic-cover}
设 \(R(y_n)\) 为 Lee--Yang 塔第 \(n\) 层的分支集。则有严格分解
\[
R(y_n)=\bigsqcup_{i=0}^{4}\bigl(P_i+E_{\min}[2^n]\bigr),
\qquad
\deg R(y_n)=5\cdot 4^n.
\]
并且在层间映射 \([2]\) 下，每个分片都按最高位截断相容传递，从而逆极限满足
\[
\mathcal R_\infty:=\varprojlim_n R(y_n)
\cong
\bigsqcup_{i=0}^{4}T_2(E_{\min})
\cong
\bigsqcup_{i=0}^{4}\ZZ_2^2.
\]
换言之，Lee--Yang 分支塔的极限对象并不是单一 \(2\)-进地址面，而是其五片不相交拷贝的有限覆盖；连续基底由 \(T_2(E_{\min})\) 承载，而分支多重性则被五值 sheet 变量精确保留。
\end{theorem}
\begin{proof}
五分片分解与分支度公式是定理 \ref{thm:xi-time-part62d-leyang-branch-inverse-limit-five-torsors} 及其上游的有限层地址分解的直接内容。再结合推论 \ref{cor:killo-2adic-tate-hypercube-limit} 对 \(T_2(E_{\min})\cong\ZZ_2^2\) 的识别，便得到所示逆极限分解。由于过渡映射在每个分片上都由 \([2]\) 诱导，故其坐标形式正是最高位截断。证毕。
\end{proof}

\begin{theorem}[连续基底的普适性与非完备性]\label{thm:conclusion-continuous-base-universal-incomplete}
在有限素局部化与 Lee--Yang 分支塔两条链上，连续基底都具有普适性，但都不是完备承载体：
\begin{enumerate}
  \item 对有限素局部化，所有 \(\Sigma_S\) 的连续商都统一坍缩为 \(\TT\)，但 \(\TT\) 本身完全遗忘素数支撑；只有附加全不连通核 \(P_S=\prod_{p\in S}\ZZ_p\) 后，算术类型才可恢复。
  \item 对 Lee--Yang 分支塔，所有分支片都共享同一个连续基底 \(T_2(E_{\min})\)，但忽略 sheet 标号后便丢失完整分支对象；只有附加离散标签 \(i\in\{0,\dots,4\}\) 后，逆极限全对象才可恢复。
\end{enumerate}
因此，在本文框架中，连续相位基底从来只能作为统一几何底座出现；一旦要求可恢复的算术记忆或分支多重性，就必须外接一条与之正交的离散账本或有限覆盖。
\end{theorem}
\begin{proof}
第一部分直接由定理 \ref{thm:conclusion-finite-localization-circlequotient-profinite-kernel} 给出：对任意有限 \(S,T\)，均有
\[
\Sigma_S/P_S\cong \TT\cong \Sigma_T/P_T,
\]
但 \(\Sigma_S\cong\Sigma_T\iff S=T\)。第二部分则是定理 \ref{thm:conclusion-leyang-five-sheet-dyadic-cover} 的直接重述：五片分支共享同一 \(T_2(E_{\min})\) 基底，而片标一旦遗忘，完整 branch 对象便不再可恢复。两部分合并即得结论。证毕。
\end{proof}

\begin{conjecture}[相位、账本与分支的最小三分裂原则]\label{conj:conclusion-phase-ledger-branch-minimal-trisplitting}
任何与扫描--投影生成器兼容、并同时忠实承载乘法账本与 Lee--Yang 分支数据的实现，都必须在最小层级上分裂为三种彼此不可替代的承载：
\[
\text{连续相位商 } \TT
\quad+\quad
\text{全不连通 prime-ledger 核 } P_S
\quad+\quad
\text{有限 sheet 覆盖 } F.
\]
更精确地，应存在一个最小因子化
\[
\mathcal X\longrightarrow \TT
\]
使其纤维再分解为
\[
P_S\times F,
\qquad
P_S=\prod_{p\in S}\ZZ_p,
\]
其中 \(F\) 为有限离散分支集；对当前 Lee--Yang 链，预期 \(|F|=5\)。删去这三部分中的任意一部分，都将分别破坏相位输运、一致的素支撑恢复或分支全对象恢复。
\end{conjecture}

\noindent
上述链条把实现层单射压缩、有限素局部化的单圆连续商、Gödel 半直积的 \(2\)-进仿射表示以及 Lee--Yang 分支塔的五片 dyadic 覆盖统一到同一结构判断：连续部分只能以商的形式出现；凡可恢复的算术记忆与分支多重性，都必须沿全不连通核或有限覆盖外置。
